\section{Main Results on 8iVFB}
\label{sec:main-results}

Figure~\ref{fig:rd-curve-main} shows the rate-distortion performance of the
proposed method compared to the PCGCv2 baseline (and Unicorn for reference),
while Figure~\ref{fig:per-rate-improvement} summarises the average PSNR
improvement across all sequences and rates.

\begin{figure}[t]
  \centering
  \includegraphics[width=0.98\textwidth]{rd_curve.png}
  \caption{%
    Rate-distortion curves for D1 PSNR (left), Edge PSNR (centre), and
    Flat PSNR (right) at the three training rates (r2, r4, r6), averaged
    over all four 8iVFB sequences.
    PCGCv2 baseline (blue) vs.\ proposed post-processed output (orange).
    Unicorn (green, dashed) serves as a reference for a higher-performing codec.
    Zero-shot transfer to Unicorn (purple, dashed) is included for completeness.
    PSNR improvements at each rate are annotated.
  }
  \label{fig:rd-curve-main}
\end{figure}

\begin{figure}[t]
  \centering
  \includegraphics[width=0.78\textwidth]{per_rate_improvement.png}
  \caption{%
    Average PSNR improvement $\delta$ (dB) over PCGCv2 baseline, broken down
    by rate and metric type (edge, flat, D1).
    Improvements are largest at r2 (0.05~bpp), where oversmoothing is most
    severe, and diminish at higher rates, consistent with the decreasing
    displacement signal.
  }
  \label{fig:per-rate-improvement}
\end{figure}

\begin{figure}[t]
  \centering
  \includegraphics[width=0.78\textwidth]{per_sequence_r2.png}
  \caption{%
    Per-sequence PSNR improvement at r2 (0.05~bpp).
    The model achieves consistently positive improvements across all four
    sequences (edge $\delta$: 1.97--2.29~dB, flat $\delta$: 1.22--1.47~dB),
    with three of the four sequences not seen during training
    (\textit{redandblack} is the held-out validation sequence).
  }
  \label{fig:per-sequence-r2}
\end{figure}

\subsection{Per-Sequence Results}

Tables~\ref{tab:results-edge-r2}--\ref{tab:results-flat} report the full
curvature-stratified results for all four sequences, three rates, and
400 frames.

\begin{table}[t]
  \centering
  \caption{%
    Edge-stratum PSNR improvement $\delta$ (dB) over PCGCv2 baseline.
    All values are average over 100 frames per sequence.
  }
  \label{tab:results-edge-r2}
  \begin{tabular}{lrrr}
    \toprule
    Sequence & r2 & r4 & r6 \\
    \midrule
    longdress   & $+2.21$ & $+0.58$ & $+0.27$ \\
    loot        & $+1.78$ & $+0.52$ & $+0.24$ \\
    redandblack & $+2.06$ & $+0.62$ & $+0.28$ \\
    soldier     & $+2.20$ & $+0.61$ & $+0.28$ \\
    \midrule
    \textbf{Average} & $\mathbf{+2.06}$ & $\mathbf{+0.58}$ & $\mathbf{+0.27}$ \\
    \bottomrule
  \end{tabular}
\end{table}

\begin{table}[t]
  \centering
  \caption{%
    Flat-stratum PSNR improvement $\delta$ (dB) over PCGCv2 baseline.
  }
  \label{tab:results-flat}
  \begin{tabular}{lrrr}
    \toprule
    Sequence & r2 & r4 & r6 \\
    \midrule
    longdress   & $+1.69$ & $+0.59$ & $+0.16$ \\
    loot        & $+1.25$ & $+0.67$ & $+0.20$ \\
    redandblack & $+1.38$ & $+0.61$ & $+0.17$ \\
    soldier     & $+1.54$ & $+0.61$ & $+0.18$ \\
    \midrule
    \textbf{Average} & $\mathbf{+1.47}$ & $\mathbf{+0.62}$ & $\mathbf{+0.18}$ \\
    \bottomrule
  \end{tabular}
\end{table}

\begin{table}[t]
  \centering
  \caption{%
    Aggregate D1 PSNR improvement $\delta$ (dB) over PCGCv2 baseline.
  }
  \label{tab:results-d1}
  \begin{tabular}{lrrr}
    \toprule
    Sequence & r2 & r4 & r6 \\
    \midrule
    longdress   & $+1.92$ & $+0.54$ & $+0.21$ \\
    loot        & $+1.46$ & $+0.54$ & $+0.21$ \\
    redandblack & $+1.76$ & $+0.56$ & $+0.23$ \\
    soldier     & $+1.85$ & $+0.58$ & $+0.23$ \\
    \midrule
    \textbf{Average} & $\mathbf{+1.75}$ & $\mathbf{+0.55}$ & $\mathbf{+0.22}$ \\
    \bottomrule
  \end{tabular}
\end{table}

\subsection{Absolute PSNR Values}

Table~\ref{tab:absolute-psnr-results} provides absolute PSNR values for both the
PCGCv2 baseline and the refined output, giving context for the relative
improvements.

\begin{table}[t]
  \centering
  \caption{%
    Absolute curvature-stratified PSNR (dB) for PCGCv2 baseline (Dec.) and
    post-processed output (Ref.), averaged over 4 sequences.
  }
  \label{tab:absolute-psnr-results}
  \small
  \begin{tabular}{lccccccc}
    \toprule
    & \multicolumn{3}{c}{Edge PSNR (dB)} & \multicolumn{3}{c}{Flat PSNR (dB)} \\
    \cmidrule(lr){2-4}\cmidrule(lr){5-7}
    & r2 & r4 & r6 & r2 & r4 & r6 \\
    \midrule
    PCGCv2 decoded & 60.05 & 65.64 & 68.17 & 62.70 & 67.95 & 70.95 \\
    Proposed (refined) & 62.11 & 66.22 & 68.43 & 64.17 & 68.57 & 71.13 \\
    \midrule
    $\Delta$ & $+2.06$ & $+0.58$ & $+0.27$ & $+1.47$ & $+0.62$ & $+0.18$ \\
    \bottomrule
  \end{tabular}
\end{table}

At r2, the refined edge PSNR of 62.11~dB approaches the baseline flat PSNR
of 62.70~dB --- the model nearly closes the stratification gap at r2
for the \textit{loot} and \textit{soldier} sequences.

\subsection{BD-PSNR Summary}

Table~\ref{tab:bd-psnr} summarises the Bjontegaard Delta metrics computed
over the evaluated rate range (r2--r6).

\begin{table}[t]
  \centering
  \caption{%
    Bjontegaard Delta PSNR (dB) over the r2--r6 rate range.
    Positive values indicate quality gain at equal rate.
  }
  \label{tab:bd-psnr}
  \begin{tabular}{lccc}
    \toprule
    Metric & BD-PSNR (dB) \\
    \midrule
    D1 PSNR    & $+0.85$ \\
    Edge PSNR  & $+0.97$ \\
    Flat PSNR  & $+0.79$ \\
    \bottomrule
  \end{tabular}
\end{table}

The BD-PSNR of $+0.97$~dB on the edge metric, integrated across rates r2--r6,
represents a substantial and consistent gain.
For reference, the improvement in G-PCC from version 11 to version 14
(across approximately two years of standardisation effort) is approximately
1--2~dB BD-PSNR on the MPEG evaluation sequences.

\subsection{Rate-Distortion Analysis}

Table~\ref{tab:rd-bpp} provides the bits-per-point values at each evaluated
rate point, providing the full rate-distortion context.

\begin{table}[t]
  \centering
  \caption{%
    PCGCv2 bits-per-point (bpp) at each evaluated rate, averaged over
    sequences and 100 frames. The bpp varies slightly per sequence and
    frame due to the content-adaptive entropy model.
  }
  \label{tab:rd-bpp}
  \begin{tabular}{lccc}
    \toprule
    Sequence & r2 (bpp) & r4 (bpp) & r6 (bpp) \\
    \midrule
    longdress   & 0.048 & 0.153 & 0.317 \\
    loot        & 0.047 & 0.149 & 0.311 \\
    redandblack & 0.050 & 0.158 & 0.325 \\
    soldier     & 0.049 & 0.155 & 0.321 \\
    \midrule
    Average     & 0.049 & 0.154 & 0.319 \\
    \bottomrule
  \end{tabular}
\end{table}

\subsection{Full Oversmoothing Characterisation}

Table~\ref{tab:full-oversmoothing} extends the oversmoothing analysis of
Chapter~\ref{ch:oversmoothing} to all seven PCGCv2 rate points, providing
the complete picture of how the stratification gap evolves with rate.

\begin{table}[t]
  \centering
  \caption{%
    Curvature-stratified reverse PSNR for PCGCv2 at all 7 rate points.
    Values are for a single representative frame per sequence.
    $\Delta = \text{Flat} - \text{Edge}$ (dB); positive = oversmoothing.
  }
  \label{tab:full-oversmoothing}
  \small
  \setlength{\tabcolsep}{4pt}
  \begin{tabular}{llrrrrr}
    \toprule
    Sequence & Rate & bpp & Flat & Medium & Edge & $\Delta$ \\
    \midrule
    \multirow{7}{*}{longdress}
      & r1 & 0.025 & 56.17 & 57.48 & 55.05 & 1.12 \\
      & r2 & 0.048 & 62.55 & 62.61 & 59.88 & 2.67 \\
      & r3 & 0.092 & 66.04 & 65.96 & 63.83 & 2.21 \\
      & r4 & 0.153 & 68.16 & 67.81 & 65.86 & 2.31 \\
      & r5 & 0.247 & 69.88 & 69.53 & 67.49 & 2.39 \\
      & r6 & 0.317 & 70.95 & 70.20 & 68.12 & 2.83 \\
      & r7 & 0.401 & 71.91 & 71.04 & 68.77 & 3.15 \\
    \midrule
    \multirow{7}{*}{loot}
      & r1 & 0.024 & 56.35 & 58.60 & 56.11 & 0.24 \\
      & r2 & 0.047 & 63.36 & 63.55 & 60.77 & 2.58 \\
      & r3 & 0.089 & 66.29 & 66.13 & 64.10 & 2.19 \\
      & r4 & 0.149 & 68.33 & 68.04 & 66.21 & 2.12 \\
      & r5 & 0.242 & 70.01 & 69.69 & 67.96 & 2.05 \\
      & r6 & 0.311 & 71.02 & 70.47 & 68.69 & 2.33 \\
      & r7 & 0.397 & 71.92 & 71.22 & 69.32 & 2.60 \\
    \midrule
    \multirow{7}{*}{redandblack}
      & r1 & 0.027 & 57.29 & 58.24 & 54.87 & 2.41 \\
      & r2 & 0.050 & 62.94 & 62.92 & 58.69 & 4.24 \\
      & r3 & 0.095 & 65.85 & 65.73 & 62.17 & 3.68 \\
      & r4 & 0.158 & 67.97 & 67.50 & 64.19 & 3.78 \\
      & r5 & 0.256 & 70.13 & 69.24 & 66.03 & 4.11 \\
      & r6 & 0.325 & 70.89 & 69.96 & 66.67 & 4.22 \\
      & r7 & 0.410 & 71.78 & 70.79 & 67.25 & 4.53 \\
    \midrule
    \multirow{7}{*}{soldier}
      & r1 & 0.025 & 57.19 & 58.16 & 55.46 & 1.73 \\
      & r2 & 0.049 & 62.51 & 62.65 & 59.88 & 2.63 \\
      & r3 & 0.093 & 65.61 & 65.45 & 63.54 & 2.07 \\
      & r4 & 0.155 & 67.40 & 67.35 & 65.76 & 1.64 \\
      & r5 & 0.252 & 69.72 & 69.25 & 67.78 & 1.95 \\
      & r6 & 0.321 & 70.73 & 70.02 & 68.47 & 2.26 \\
      & r7 & 0.407 & 71.74 & 70.82 & 69.19 & 2.55 \\
    \bottomrule
  \end{tabular}
\end{table}

\subsection{G-PCC Comparison: Full Oversmoothing Table}

Table~\ref{tab:gpcc-oversmoothing} provides the complete G-PCC oversmoothing
data for comparison.

\begin{table}[t]
  \centering
  \caption{%
    Curvature-stratified reverse PSNR for G-PCC at its evaluated rate points
    (Trisoup mode).
    $\Delta = \text{Flat} - \text{Edge}$ (dB).
  }
  \label{tab:gpcc-oversmoothing}
  \small
  \setlength{\tabcolsep}{4pt}
  \begin{tabular}{llrrrrr}
    \toprule
    Sequence & Rate & bpp & Flat & Medium & Edge & $\Delta$ \\
    \midrule
    \multirow{5}{*}{longdress}
      & ts6 & 0.004 & 39.01 & 39.79 & 38.49 & 0.51 \\
      & ts5 & 0.013 & 52.09 & 51.74 & 48.68 & 3.41 \\
      & ts4 & 0.043 & 61.76 & 60.95 & 57.09 & 4.67 \\
      & ts3 & 0.133 & 67.23 & 65.90 & 63.98 & 3.25 \\
      & ts2 & 0.383 & 68.70 & 67.47 & 66.98 & 1.72 \\
    \midrule
    \multirow{5}{*}{loot}
      & ts6 & 0.005 & 39.12 & 40.84 & 38.88 & 0.24 \\
      & ts5 & 0.014 & 51.14 & 51.74 & 46.72 & 4.42 \\
      & ts4 & 0.043 & 63.08 & 62.27 & 55.26 & 7.82 \\
      & ts3 & 0.128 & 66.00 & 65.77 & 62.68 & 3.31 \\
      & ts2 & 0.372 & 68.79 & 67.49 & 66.98 & 1.81 \\
    \midrule
    \multirow{5}{*}{redandblack}
      & ts6 & 0.005 & 39.67 & 39.70 & 38.79 & 0.88 \\
      & ts5 & 0.013 & 51.18 & 50.86 & 44.20 & 6.98 \\
      & ts4 & 0.044 & 61.78 & 60.91 & 53.90 & 7.88 \\
      & ts3 & 0.137 & 66.45 & 65.65 & 62.24 & 4.21 \\
      & ts2 & 0.403 & 68.52 & 67.35 & 66.49 & 2.03 \\
    \midrule
    \multirow{5}{*}{soldier}
      & ts6 & 0.005 & 42.54 & 42.38 & 39.18 & 3.37 \\
      & ts5 & 0.013 & 50.30 & 49.98 & 47.22 & 3.09 \\
      & ts4 & 0.044 & 62.86 & 61.34 & 57.70 & 5.16 \\
      & ts3 & 0.135 & 67.08 & 65.65 & 63.73 & 3.35 \\
      & ts2 & 0.387 & 62.61 & 62.72 & 62.68 & $-0.07$ \\
    \bottomrule
  \end{tabular}
\end{table}

The G-PCC data reveals a more variable pattern than PCGCv2.
At the lowest rate (ts6), G-PCC shows minimal stratification gap ($\Delta < 1$~dB
for most sequences) because its artifacts are approximately uniform quantisation
noise.
However, at intermediate rates (ts4--ts5), G-PCC exhibits a \emph{larger} gap
than PCGCv2 at comparable bitrates, particularly for \textit{loot}
($\Delta = 7.82$~dB) and \textit{redandblack} ($\Delta = 7.88$~dB).
This reflects the Trisoup triangle approximation artefacts, which can
introduce large errors at surface discontinuities when the triangle size is
poorly matched to the local geometry.
At high rates (ts2), G-PCC's gap drops below 2~dB as the triangles are small
enough to represent fine detail.
In contrast, PCGCv2's gap \emph{increases} monotonically with rate
(Table~\ref{tab:full-oversmoothing}).

% ============================================================
